% =============================================================================
% Autore originale: Fabio Tempera
% Licenza: MIT (vedere LICENSE nel medesimo direttorio)
%
% DOCUMENTO
%   Lettera alle ARPA regionali per proporre una rete regionale di misura della
%   brillanza del cielo notturno e per sollecitare un finanziamento regionale
%   dedicato all'installazione delle stazioni presso gli enti del territorio.
%
% COMPILAZIONE
%   pdflatex arpa.tex  (due passate)
%   Classe report, pacchetti in packages.sty. Non usare Markdown.
% =============================================================================

\documentclass[11pt,a4paper]{report}

\usepackage{packages}

\definecolor{nightblue}{RGB}{22,38,58}
\definecolor{steel}{RGB}{80,100,120}
\definecolor{linegray}{RGB}{200,200,200}

\pagestyle{fancy}
\fancyhf{}
\renewcommand{\headrulewidth}{0.4pt}
\renewcommand{\headrule}{\hbox to\headwidth{\color{linegray}\leaders\hrule height \headrulewidth\hfill}}
\fancyhead[L]{\small\color{steel} CieloBuio OdV}
\fancyhead[R]{\small\color{steel} Proposta alle ARPA regionali, rete TESS-4C}
\fancyfoot[C]{\small\color{steel}\thepage}

\hypersetup{
  colorlinks=true,
  urlcolor=nightblue,
  pdftitle={Proposta di collaborazione alle ARPA regionali. Rete TESS-4C},
  pdfauthor={Fabio Tempera}
}

\setlist[itemize]{leftmargin=1.4em, itemsep=0.25em, topsep=0.3em}

\begin{document}

\noindent Alla Direzione generale

\noindent dell'Agenzia regionale per la protezione dell'ambiente

\vspace{0.8cm}

\noindent \textit{Oggetto}: proposta di collaborazione per una rete regionale di misura della brillanza del cielo notturno

\vspace{0.8cm}

\noindent Con questa lettera si propone un'ipotesi di collaborazione per estendere su base regionale la rete TESS-4C, promossa da CieloBuio OdV.

La rete è costituita da stazioni fotometriche a terra che misurano la brillanza del cielo notturno in modo continuo e confrontabile, collegate alla rete europea STARS4ALL. Lo strumento è il fotometro TESS-4C, a 4 bande, alimentato a 5\,V e dotato di trasmissione dati via WiFi.

Il monitoraggio dell'efficacia delle norme regionali sul contenimento della luce artificiale richiede dati a terra, puntuali e continui. Le stime satellitari restano utili per inquadrare il territorio, ma non sostituiscono una misura locale.

In questo quadro si propone quanto segue.

\begin{itemize}
\item Una valutazione da parte dell'Agenzia dell'interesse a una rete regionale, con l'eventuale inserimento di stazioni presso enti del territorio.
\item La richiesta di fondi regionali dedicati, a sostegno dell'acquisto e dell'installazione degli strumenti presso gli enti ospitanti.
\item Un incontro tecnico per definire siti, requisiti di installazione e modalità di calibrazione.
\end{itemize}

Il finanziamento regionale permetterebbe di sostenere l'installazione senza che il progetto dipenda da bandi di altra natura. I dati delle stazioni resterebbero consultabili in formato aperto, con una sintesi pubblica per territorio.

Per approfondimenti o per fissare un incontro è possibile contattare il referente del progetto ai recapiti seguenti.

\vspace{0.6cm}

\noindent Cordiali saluti.

\vspace{1.2cm}

\noindent Fabio Tempera, CieloBuio OdV

\noindent \url{fabiot2801@gmail.com} \quad$\cdot$\quad WhatsApp 329 883 5142

\end{document}
